% =========================================================
% Defining the Order on Q
% =========================================================

\label{sec:defining-order-on-q}

\begin{exposition}
The rational numbers are not merely a field. They also carry an order
relation compatible with the algebraic operations, so that $\mathbb{Q}$
becomes an ordered field.
\end{exposition}

\begin{remark*}
For many arguments it is cleaner to encode order using the set of
positive elements rather than to take $<$ as primitive.
\end{remark*}

% Positive cone now defined live in arithmetic-operations-relations/algebraic-structures (def:positive-cone).

\begin{remark*}
This viewpoint makes the interaction between algebra and order
transparent. Inequalities can be translated into statements about
positive elements and then handled algebraically.
\end{remark*}

\begin{definition}[Ordered integral domain]
\label{def:ordered-integral-domain}
An \emph{ordered integral domain} is an integral domain
\[
(F,+,\cdot,0,1)
\]
together with a relation $<$ on $F$ such that:
\begin{enumerate}
  \item for all $a,b\in F$, exactly one of the relations
  \[
  a<b,\qquad a=b,\qquad b<a
  \]
  holds;
  \item if $a<b$, then for every $c\in F$,
  \[
  a+c<b+c;
  \]
  \item if $a<b$ and $0<c$, then
  \[
  ac<bc.
  \]
\end{enumerate}
\end{definition}

\begin{definition}[Ordered field]
\label{def:rational-ordered-field}
An \emph{ordered field} is a field
\[
(F,+,\cdot,0,1)
\]
together with a relation $<$ on $F$ making $F$ into an ordered integral
domain.
\end{definition}

\begin{definition}[Ordered subfield]
\label{def:ordered-subfield}
Let $E$ and $F$ be ordered fields with $E\subseteq F$. We say that $E$ is an
\emph{ordered subfield} of $F$ if $E$ is a subfield of $F$ and the order on
$E$ is the restriction of the order on $F$.
\end{definition}

\begin{definition}[Densely ordered system]
\label{def:densely-ordered-system}
Let $S$ be a set with a relation $<$. We say that $(S,<)$ is
\emph{densely ordered} if:
\begin{enumerate}
  \item $(S,<)$ is a simply ordered system; and
  \item for every $x,y\in S$ with $x<y$, there exists $z\in S$ such that
  \[
  x<z<y.
  \]
\end{enumerate}
\end{definition}
